% ============================================================
\section{Test LLM}
\label{sec:llm}
% ============================================================

\subsection{BrokerImpl}

Il codice tra \texttt{\#\#\#Test START\#\#} e \texttt{\#\#\#Test END\#\#} viene estratto dalla
risposta e copiato in
\texttt{openjpa-kernel/src/test/java/.../kernel/}. Dove un test rivela un'assunzione errata dell'LLM, il metodo è marcato
\texttt{@Ignore}. Dove invece un test documenta un fault reale,
l'asserzione riflette il comportamento osservato (quello col fault), non quello atteso secondo
la documentazione: il test resta attivo e passa.

In quasi tutti i prompt il conteggio di test dichiarato dall'LLM nel testo non coincide con
quello reale nel codice generato.

\subsubsection{Z1 (\hyperref[prompt:z1]{Appendice}).}
Sceglie 10 metodi evitando di propria iniziativa tutti quelli che richiedono
\texttt{initialize()}, la stessa precondizione già nota (Capitolo~\ref{sec:implementazione}) e
riscontrata da Randoop per errore a runtime, qui evitata a monte per lettura del codice. 21 test
generati (uno oltre il tetto di 20). 20 passano, 1 escluso con \texttt{@Ignore}:
\texttt{testSetRollbackOnlyWithoutActiveTransactionThrows} si aspettava
\texttt{NoTransactionException}, il comportamento reale lancia \texttt{InvalidStateException}.

\subsubsection{Z2 (\hyperref[prompt:z2]{Appendice}).}
Sottoinsieme di metodi diverso da Z1. Usa la reflection una sola volta per forzare
\texttt{\_initializeWasInvoked} e testare \texttt{clone()}, giustificandolo come
``arrangiare la precondizione'' e non ``asserire su stato interno''. È una distinzione discutibile
rispetto al vincolo black-box del prompt, dato che manipolare un campo privato resta comunque una
tecnica white-box. 20 test, entro il tetto. 18 passano, 2 esclusi con \texttt{@Ignore}: si
aspettava \texttt{null} da \texttt{getConnectionFactoryName()}/\texttt{getConnectionFactory2Name()}
su un broker nuovo; il default reale è una stringa vuota (\texttt{\_connectionFactoryName = ""}
nel sorgente).

\subsubsection{F1 (\hyperref[prompt:f1]{Appendice}).}
Terzo insieme di metodi, incluso l'overload senza argomenti \texttt{lock()}/\texttt{unlock()}. 20 test, entro
il tetto. Compila senza errori, tutti passano. A differenza di Z2, modella
correttamente sia \texttt{getConnectionFactoryName()} sia \texttt{getConnectionFactory()}.

\subsubsection{F2 (\hyperref[prompt:f2]{Appendice}).}
Copre 10 metodi diversi dai precedenti, incluso \texttt{getCacheFinderQuery}/
\texttt{setCacheFinderQuery}. 20 test, entro il tetto. Compila senza errori e tutti
passano, incluso un test che documenta un \textbf{fault reale} verificato sul sorgente:
\texttt{setCacheFinderQuery(boolean)} scrive per errore sul campo \texttt{\_cachePreparedQuery}
invece che su \texttt{\_cacheFinderQuery}, che risulta quindi mai modificabile dal proprio setter
pubblico.

\subsubsection{C1 (\hyperref[prompt:c1]{Appendice}).}
Il ragionamento richiesto dal prompt è visibile nel testo: per ciascun metodo scelto l'LLM motiva
la selezione prima di scrivere codice. 19 test, entro il tetto. Compila senza errori e
tutti passano. Ritrova indipendentemente lo stesso fault di F2 su \texttt{setCacheFinderQuery}, e
ne trova un secondo, non emerso altrove: la Javadoc di \texttt{getSuppressBatchOLELogging()}
dichiara ``Defaults to true'', ma il campo \texttt{\_suppressBatchOLELogging} è dichiarato
\texttt{false}.

\subsubsection{C2 (\hyperref[prompt:c2]{Appendice}).}
Stesso schema di C1, con l'esempio few-shot come riferimento di stile; insieme di metodi quasi
identico. 19 test, entro il tetto. Compila senza errori e tutti passano. Ritrova lo
stesso fault su \texttt{setCacheFinderQuery} per la terza volta, qui
verificato nel modo più rigoroso: legge con la reflection entrambi i campi privati dopo la
chiamata, confermando che \texttt{\_cacheFinderQuery} resta \texttt{true} mentre
\texttt{\_cachePreparedQuery} diventa \texttt{false}.

\subsection{LoginDialog}
% non affrontato in questo giro, per decisione presa in Sezione~\ref{sec:randoop}
